%% ma: L12-B08-0153
%% lop: 12
%% chuong: 2
%% bai: 8
%% dang: 12.8.51
%% loai: TN4
%% muc: TH
%% dap_an: "A"
%% nguon: "CĐ2.4.pdf giáo viên gửi 27/10/2026 (M07, Lê Minh Khôi – Nguyễn Anh Khôi, Chương II Chủ đề 4 Các bài toán về tâm tỉ cự), Câu 3"
%% kien_thuc: "Bài 8 (biểu thức tọa độ của phép toán vectơ, độ dài vectơ), Bài 6 (điểm thỏa hệ thức vectơ, quy tắc ba điểm)"
%% trang_thai: cho-duyet
%% ly_do: "Dạng mới 12.8.51 đề xuất, chờ giáo viên duyệt"
\begin{ex}
Trong không gian $Oxyz$, cho ba điểm $A(4;2;2)$, $B(1;1;-1)$, $C(2;-2;-2)$. Tọa độ của điểm $M$ thỏa mãn $\left|\overrightarrow{MA}+2\overrightarrow{MB}-\overrightarrow{MC}\right|$ đạt giá trị nhỏ nhất là
\choice{\True $M(2;3;1)$}{$M(0;3;1)$}{$M(0;-3;1)$}{$M(0;1;2)$}
\loigiai{Gọi $I$ thỏa $\overrightarrow{IA}+2\overrightarrow{IB}-\overrightarrow{IC}=\vec 0\Rightarrow I=\dfrac{A+2B-C}2=(2;3;1)$.

$\left|\overrightarrow{MA}+2\overrightarrow{MB}-\overrightarrow{MC}\right|=2MI$ nhỏ nhất bằng $0$ khi $M\equiv I(2;3;1)$.}
\end{ex}
